\subsection{Fake-data tests of the estimator's bias}
\label{sec:fived16-fd}
The fake-data tests compare the estimator's bias at truths other than the simulation with the reported covariance. Pseudo-data are generated at \fvval{f0b-all-truths} truths and unfolded with the \estMultiD{}. The mean unfolded result is compared with each truth on the 16 cells, through $\chi^2$ with 16 degrees of freedom and the per-cell pull $z$. The truths are the simulation's prior and the following. The five hybrid truths of the generator tests (Sec.~\ref{sec:joint5d}) are each a prediction on the fine grid. The historical shapes of \S\ref{sec:gbdt-model-dependence} are W1, W2 and W3. The remaining truths are a GiBUU-to-GENIE $\Eavail$ shape, a noise-free closure at the NuWro shape and a synthetic tilt of the $q_3$ distribution at fixed $(\Eavail,W)$. A truth from a $C_{\mathrm{model}}$ generator is compared with the covariance without that generator's shift. A truth is listed as failing when $p<\fvval{trig-p}$ or when any cell has $|z|>\fvval{trig-z}$ (pre-specified; Table~\ref{tab:test-index}). These are not held-out truths, and the detector response is not probed. Table~\ref{tab:fived16-fd} gives each truth under the reported covariance, without $C_{\mathrm{PL,5D}}$, and under $C_e^{\mathrm{infl}}+C_{\mathrm{model}}$ alone. The bias of each truth is the same in all three; only the covariance differs.
% - Referee round 20261104, PRD read 1, R3 (a): FD-DIRECTOR-RULINGS.md sec. 2 (092702292, 2026-10-25 10:32 -0700; sha256 2bfba459806310a3), before F0's declaration (7d7dac1d3) and report (28c4870f4).

\begin{table}[htbp]
\centering\footnotesize
\caption[Fake-data tests of the sixteen-cell bias]{Under each of the three covariances, the fake-data bias meets the listing criterion at the GiBUU hybrid truth only. Under the baseline covariance, $C_{\mathrm{rep}}$ with $C_e$ in place of $C_e^{\mathrm{infl}}$, the $q_3$ tilt was also listed, and the two smaller covariances listed more truths, with the same bias. Fake-data tests on the sixteen cells: $\chi^2$ (16 degrees of freedom), $p$ and the cells with $|z|>\fvval{trig-z}$ of the bias of the mean against the reported $C_{\mathrm{rep}}$, against $C_e^{\mathrm{infl}}+C_{\mathrm{model}}+C_{\mathrm{resp}}$ (without the run-period block), and against $C_e^{\mathrm{infl}}+C_{\mathrm{model}}$ alone, each with the truth's own generator shift removed where the second column says so ($-$). $n$ is the number of pseudo-experiments.}
\label{tab:fived16-fd}
\scriptsize\setlength{\tabcolsep}{3pt} \resizebox{\textwidth}{!}{
\begin{tabular}{@{}llrrrlrrlrrl@{}}
\toprule
truth & covariance & $n$ & \multicolumn{3}{c}{reported} & \multicolumn{3}{c}{without $C_{\mathrm{PL,5D}}$} & \multicolumn{3}{c}{$C_{\mathrm{total}}$} \\
 & & & $\chi^2$ & $p$ & cells & $\chi^2$ & $p$ & cells & $\chi^2$ & $p$ & cells \\
\midrule
simulation prior (MINERvA Tune~v1) & full & \fvval{f0b-nominal-n} & \fvval{f0b-nominal-chi2} & $\fvval{f0b-nominal-p}$ & \fvval{f0b-nominal-cells} & \fvval{f0b-nominal-chi2-nopl} & $\fvval{f0b-nominal-p-nopl}$ & \fvval{f0b-nominal-cells-nopl} & \fvval{f0-nominal-chi2} & $\fvval{f0-nominal-p}$ & \fvval{f0-nominal-cells} \\
MINERvA Tune~v1, hybrid & full & \fvval{f0b-A7_mnvtune-n} & \fvval{f0b-A7_mnvtune-chi2} & $\fvval{f0b-A7_mnvtune-p}$ & \fvval{f0b-A7_mnvtune-cells} & \fvval{f0b-A7_mnvtune-chi2-nopl} & $\fvval{f0b-A7_mnvtune-p-nopl}$ & \fvval{f0b-A7_mnvtune-cells-nopl} & \fvval{f0-A7_mnvtune-chi2} & $\fvval{f0-A7_mnvtune-p}$ & \fvval{f0-A7_mnvtune-cells} \\
GENIE CV, hybrid & full & \fvval{f0b-A7_genie_cv-n} & \fvval{f0b-A7_genie_cv-chi2} & $\fvval{f0b-A7_genie_cv-p}$ & \fvval{f0b-A7_genie_cv-cells} & \fvval{f0b-A7_genie_cv-chi2-nopl} & $\fvval{f0b-A7_genie_cv-p-nopl}$ & \fvval{f0b-A7_genie_cv-cells-nopl} & \fvval{f0-A7_genie_cv-chi2} & $\fvval{f0-A7_genie_cv-p}$ & \fvval{f0-A7_genie_cv-cells} \\
GENIE + MEC, hybrid & $-$GENIE + MEC & \fvval{f0b-A7_genie_mec-n} & \fvval{f0b-A7_genie_mec-chi2} & $\fvval{f0b-A7_genie_mec-p}$ & \fvval{f0b-A7_genie_mec-cells} & \fvval{f0b-A7_genie_mec-chi2-nopl} & $\fvval{f0b-A7_genie_mec-p-nopl}$ & \fvval{f0b-A7_genie_mec-cells-nopl} & \fvval{f0-A7_genie_mec-chi2} & $\fvval{f0-A7_genie_mec-p}$ & \fvval{f0-A7_genie_mec-cells} \\
GENIE + MEC, $(\Eavail,W)$ shape (W2) & $-$GENIE + MEC & \fvval{f0b-W2-n} & \fvval{f0b-W2-chi2} & $\fvval{f0b-W2-p}$ & \fvval{f0b-W2-cells} & \fvval{f0b-W2-chi2-nopl} & $\fvval{f0b-W2-p-nopl}$ & \fvval{f0b-W2-cells-nopl} & \fvval{f0-W2-chi2} & $\fvval{f0-W2-p}$ & \fvval{f0-W2-cells} \\
NuWro, hybrid & $-$NuWro & \fvval{f0b-A7_nuwro-n} & \fvval{f0b-A7_nuwro-chi2} & $\fvval{f0b-A7_nuwro-p}$ & \fvval{f0b-A7_nuwro-cells} & \fvval{f0b-A7_nuwro-chi2-nopl} & $\fvval{f0b-A7_nuwro-p-nopl}$ & \fvval{f0b-A7_nuwro-cells-nopl} & \fvval{f0-A7_nuwro-chi2} & $\fvval{f0-A7_nuwro-p}$ & \fvval{f0-A7_nuwro-cells} \\
NuWro, $\Eavail$ shape (W1) & $-$NuWro & \fvval{f0b-W1-n} & \fvval{f0b-W1-chi2} & $\fvval{f0b-W1-p}$ & \fvval{f0b-W1-cells} & \fvval{f0b-W1-chi2-nopl} & $\fvval{f0b-W1-p-nopl}$ & \fvval{f0b-W1-cells-nopl} & \fvval{f0-W1-chi2} & $\fvval{f0-W1-p}$ & \fvval{f0-W1-cells} \\
NuWro, $(\pT,\pPar,\Eavail)$ shape (W3) & $-$NuWro & \fvval{f0b-W3-n} & \fvval{f0b-W3-chi2} & $\fvval{f0b-W3-p}$ & \fvval{f0b-W3-cells} & \fvval{f0b-W3-chi2-nopl} & $\fvval{f0b-W3-p-nopl}$ & \fvval{f0b-W3-cells-nopl} & \fvval{f0-W3-chi2} & $\fvval{f0-W3-p}$ & \fvval{f0-W3-cells} \\
NuWro, closure (noise-free) & $-$NuWro & \fvval{f0b-closure_nuwro-n} & \fvval{f0b-closure_nuwro-chi2} & $\fvval{f0b-closure_nuwro-p}$ & \fvval{f0b-closure_nuwro-cells} & \fvval{f0b-closure_nuwro-chi2-nopl} & $\fvval{f0b-closure_nuwro-p-nopl}$ & \fvval{f0b-closure_nuwro-cells-nopl} & \fvval{f0-closure_nuwro-chi2} & $\fvval{f0-closure_nuwro-p}$ & \fvval{f0-closure_nuwro-cells} \\
GiBUU, $\Eavail$ shape & $-$GiBUU & \fvval{f0b-stage0c_gibuu-n} & \fvval{f0b-stage0c_gibuu-chi2} & $\fvval{f0b-stage0c_gibuu-p}$ & \fvval{f0b-stage0c_gibuu-cells} & \fvval{f0b-stage0c_gibuu-chi2-nopl} & $\fvval{f0b-stage0c_gibuu-p-nopl}$ & \fvval{f0b-stage0c_gibuu-cells-nopl} & \fvval{f0-stage0c_gibuu-chi2} & $\fvval{f0-stage0c_gibuu-p}$ & \fvval{f0-stage0c_gibuu-cells} \\
GiBUU, hybrid & $-$GiBUU & \fvval{f0b-A7_gibuu-n} & \fvval{f0b-A7_gibuu-chi2} & $\fvval{f0b-A7_gibuu-p}$ & \fvval{f0b-A7_gibuu-cells} & \fvval{f0b-A7_gibuu-chi2-nopl} & $\fvval{f0b-A7_gibuu-p-nopl}$ & \fvval{f0b-A7_gibuu-cells-nopl} & \fvval{f0-A7_gibuu-chi2} & $\fvval{f0-A7_gibuu-p}$ & \fvval{f0-A7_gibuu-cells} \\
synthetic $q_3$ tilt & full & \fvval{f0b-q3-n} & \fvval{f0b-q3-chi2} & $\fvval{f0b-q3-p}$ & \fvval{f0b-q3-cells} & \fvval{f0b-q3-chi2-nopl} & $\fvval{f0b-q3-p-nopl}$ & \fvval{f0b-q3-cells-nopl} & \fvval{f0-q3-chi2} & $\fvval{f0-q3-p}$ & \fvval{f0-q3-cells} \\
\bottomrule
\end{tabular}
}
\end{table}

Under the reported covariance, only the GiBUU hybrid truth meets the listing criterion ($p<\fvval{trig-p}$, or a cell with $|z|>\fvval{trig-z}$), by $p$ and by its cells \fvval{f0b-A7_gibuu-cells}. The $q_3$ tilt, listed by $p$ alone under the baseline covariance, is not listed under $C_{\mathrm{rep}}$ ($p=\fvval{f0b-q3-p}$). Without $C_{\mathrm{PL,5D}}$, the same single truth meets it: the GiBUU hybrid truth, by its cells \fvval{f0b-A7_gibuu-cells-nopl}. The GiBUU $\Eavail$ shape, listed without $C_{\mathrm{PL,5D}}$ under the baseline covariance by its cell R06, is not listed without $C_{\mathrm{PL,5D}}$ under $C_{\mathrm{rep}}$ ($z=\fvval{f0b-stage0c_gibuu-zmax-nopl}$, against \fvval{f0b-stage0c_gibuu-zmax} with the block). The fake-data bias therefore exceeds the stated thresholds at the GiBUU hybrid truth alone, with or without $C_{\mathrm{PL,5D}}$. The bias at each truth is the same under both covariances, so the difference is the covariance's. It is not a smaller bias or a validated uncertainty. At the other \fvval{f0b-nunlisted} truths the bias meets both criteria under the reported covariance. These listings use the reported construction of $C_{\mathrm{model}}$; two alternative constructions are in the records document (Ref.~\cite{NoteRecords2026}).
% U9 (currency check of 2fe070fec): the RSN arms G6 and SM, fd/referee-sensitivity-20261111/connect-post/run-20261112/rsn_readout.json (6291f48d) and regrade.json (cb5f5d17), commit 1a8f4b3c2, by rsn_values() in make_values_fived16.py; wording as the paper's section 5.2 at 5bbb9d269.

Three validity rules classify each cell from these tests; none is a coverage test, and no cell is excluded or changed by them. Under the ATLAS User Guide's thresholds applied per cell (V-ATLAS), \lrval{vr-atlas-nvalid} cells are valid, every cell except \lrval{vr-atlas-not}, where the GiBUU hybrid truth alone gives $z=\lrval{vr-R02-z-gibuu}$ and $\lrval{vr-R06-z-gibuu}$. The guide states no rule comparing the bias spread with the spread of the predictions, so the second rule (V-PRIOR) is ours: a cell is valid when the spread of its fake-data biases is below \lrval{vr-prior-tau} of the spread of the generator predictions. Under it \lrval{vr-prior-nvalid} cells are valid, every cell except \lrval{vr-prior-not}, and \lrval{vr-both-n} cells are valid under both rules. The third, the bias-half rule, asks that a cell's bias be at most half of the truth's departure from the prior, $|b_c|\leq0.5\,|d_c|$, with the bias $b_c$ and the departure $d_c$ both relative to the prior cross section (\S\ref{sec:fived16-prr}), at each graded truth where that departure is at least \SI{1}{\percent}. A cell is unresolved if its class changes within the noise of the estimate. The rule is a diagnostic that changes nothing in the reported result, and it was defined after the fake-data biases it classifies had been seen. Under it \fvval{vhalf-n-valid} cells pass, \fvval{vhalf-n-notvalid} fail (\fvval{vhalf-cells-notvalid}) and \fvval{vhalf-n-unresolved} are unresolved. V-ATLAS bounds the bias relative to the reported uncertainty and the bias-half rule bounds it relative to the truth's departure from the prior, so the two measure different properties of the bias.
% U2 (currency check of 2fe070fec): the bias-half rule (DECISION-20261109-joseph-bias-half-rule.md) from the V-HALF readout, docs/orchestration/state/2d-closeout/valid-region-20261107/v-half/receipts/read/readout.json (e7423390; commit 35577ba6e), by vhalf_verdicts() in make_values_fived16.py; wording as the paper's section 5.2 at 5bbb9d269. Its timing (defined after the fake-data biases had been seen): DECISION-20261109-joseph-bias-half-rule.md (Joseph's decision of 2026-11-09, after F0b's grading) and the paper's Supplement at 5bbb9d269, main_supplement.tex:116 ("a diagnostic that changes nothing in the reported result and was defined after the fake-data biases it classifies had been seen").
% Note round 1 (2026-11-08, draft for a PRD-referee read and Joseph): OUTCOME-20261107-valid-region.md @ b551b05f, CITABLE
% FOR :21-25 (each cell's V-ATLAS and V-PRIOR class; the record that the guide states no prior-spread rule); NOT CITABLE
% FOR :27-37 respected (no coverage, no measurement, no cause, V-PRIOR not ATLAS's or the advisor's, a non-valid cell not
% excluded, no localization, reading (3) not used). V-PRIOR is to be amended if a source for the metric is found
% (BINDING-1 ruling 5). Numbers: \lrval{vr-*} from readout.json 91bb604e.

The truths form three groups, evaluated separately, not pooled: the 8 truths from the $C_{\mathrm{model}}$ generators; the other truths, the GENIE CV hybrid and the $q_3$ tilt; and the prior truths, the simulation prior and the Tune~v1 hybrid.
% U7 (currency check of 2fe070fec): T2B, fd/tier2-20261110/t2b/t2b.json (0e281a61; commit 8aa7392f3) and READING-T2-20261110.md (commit aaeb56f56), whose structure t2b_values() in make_values_fived16.py asserts; wording as the paper's section 5.2 at 5bbb9d269.

One control check failed: the offset check, which compares the offset corrections of two hybrid truths, gives $\chi^2=\fvval{f0b-control7-chi2}$ for 16 degrees of freedom. As fixed before any of the tests was graded, the six power alternatives built on that correction are therefore descriptive and not judged against the failure criteria; the records document gives them, with two descriptive variants of each tested truth (Ref.~\cite{NoteRecords2026}). % "eight": the eight passed controls of fv@f0b-controls-passed (values_fived16.tex:1574, "0, 1, 2, 3, 4, 5, 6, 8"); their internal numbers are left to the record (history-cuts/sec_fived16.md T10).

% Step (c), 1A (fived16/c1a/apply_c1a.py): FD-F1-REPORT.md sec. 1.1 (the grade; control 9), sec. 1.3 (the overlaps); fd_f1_readout.json (cbe411a9) labels.f1 and labels.not; FD-F1-DECLARATION.md sec. 3 (the truths as pvl_ratios shape ratios, as the C_model vertices), sec. 4 (20 Poisson experiments on half the MC; the noise-free closure), :297-298 (the CCQE and 2p2h models: NievesQELCCPXSec for both; NievesSimoVacasMECPXSec for G18, SuSAv2MECPXSec for AR23); FD-F1-BINDING.md ruling 2 (graded as non-vertex truths against C_total + C_resp; the leave-out-GENIE+MEC grade descriptive; partial evidence; not two independent tests). The tunes were generated for this test: LIMITATION-PLAN.md sec. 1.4, 1A (no GENIE v3 events existed). Not a coverage claim.
Two truths are held out of the analysis: the GENIE~v3 tunes \fvval{f1-tune-g18} and \fvval{f1-tune-ar23}, generated for this test. At each tune, \fvval{f1-n} Poisson pseudo-experiments on half the simulation and a noise-free closure are tested against the full reported covariance, as truths with no shift in $C_{\mathrm{model}}$, in a fourth group, not pooled with the others. The bias meets both criteria in all four records, the pseudo-experiment mean and the closure at each tune.

In every cell of the four records, the bias is below \fvval{f1-bsig-ceil} of the cell's reported $\sigma$; in cell R11 at \fvval{f1-tune-g18}, the bias lies below the reported uncertainty only with $C_{\mathrm{model}}$ included.

Both are GENIE tunes, so they are partial evidence and not independent tests.

The test is a bias test, not a coverage test: its Poisson pseudo-experiments do not draw the systematic parameters, so they cannot test the interval widths, and the detector response is not probed. The coverage of the result's intervals has not been validated. A first coverage test at these two tunes stopped at one of its gates before any coverage pseudo-experiment ran, and a staged redesign of it also ended without a coverage result (App.~\ref{app:devhistory}). That stop of the redesign concerns how the drawn universes were combined, not the estimator's coverage, and says nothing about whether the intervals cover, in either direction.
% absence-ok: "The coverage of the result's intervals has not been validated" is unchanged, moved up from sec_fived16.tex:283 @2fe070fec (item 15); its support is this section's three coverage tests (the first stopped at a gate, the redesign stopped at its first stage, the second descriptive with joint under-coverage), as the paper's limitation 6 at 5bbb9d269 states it
% F16-2 (2026-11-16): the tests numbered as the paper and Supplement number them (paper_body.tex "The first five-observable
% test"; supplemental_body.tex S8: "A first coverage test ... A staged redesign then stopped at its first stage"): this gate-
% stopped test is the first, the staged run below is its redesign, and LL5D is the second; App. I's two rows say the same.

A later investigation of that gate at the nominal truth, which concerns the test and neither validates nor undermines the reported interaction block, and a pre-check of the flux and detector uncertainties are in the records document.
% LL5D paragraph (2026-11-16; the coordinator's go, step 2): DECISION-20261116-joseph-paper-batch.md items 14, 21 and
% 24, mirroring the paper's approved sentences at draft/methods-short-20261109 037431430 (the Table 4 caption's definition
% of Descriptive, the section 6.3 repair paragraph, and the section 7 sentences), with "stop" in lowercase (item 21).
% LL5D-DECLARATION.md sec. 13's note items: sec. 6's resolution and size (:482-486: at most 0.025 over the cells and
% 0.025 joint per tune and variant; about 0.10 per cell, \fvval{ll5d-cellres}, and 0.08 jointly), the overlaps (sec. 3.1,
% "held out is not independent": the tune overlaps stated above), the d_T-at-means block standing in for a per-toy block
% (:288-293, :323-324), and the laterals carried (sec. 3.8, :274-280). The repair: LL5D-AMENDMENT-A1A2.md (vi keyed at
% float32; no rule, threshold, bound or cell changed, :18-26); the route repaired and the status Descriptive,
% DECISION-20261109-joseph-coverage-case-wording.md item 2(b). Values: \fvval{ll5d-*}, ported from the paper branch
% (check_ll5d_port.py); the per-tune classes are hand-typed and asserted against the readout by make_values_fived16.py
% (LL5D_PRINTED_CLASS).
\paragraph{The second coverage test.} A second coverage test at the same two held-out tunes uses the reported estimator. It has two variants, without and with the GEANT4 neutron and proton bands. For a covariance without the response and run-period terms, it flags joint under-coverage at both tunes in both variants. Without the two bands, the test at \fvval{f1-tune-g18} flags under-coverage jointly, with no cell flagged, and at \fvval{f1-tune-ar23} it flags under-coverage jointly and in \fvval{ll5d-ar23-minus2-names}. With the two bands, the flags are the same: jointly, with no cell flagged, at \fvval{f1-tune-g18}, and jointly and in \fvval{ll5d-ar23-all-names} at \fvval{f1-tune-ar23}. Over the four readings, \fvval{ll5d-njoint-min} to \fvval{ll5d-njoint-max} of \fvval{ll5d-ntoys} pseudo-experiments lie in the joint region, against a threshold of \fvval{ll5d-kjoint-all}, and the mean $\chi^2$ per degree of freedom is \fvval{ll5d-chi2ndf-min} to \fvval{ll5d-chi2ndf-max}. Whether the joint failure persists with those terms included remains untested. Under the rule fixed before the result, the test is reported as descriptive: the criteria fixed in advance were applied, but a repair was chosen after an entry-check stop (App.~\ref{app:devhistory}). Its design, the bound on the GEANT4 bands' block and the scope of its readings are in the records document's sixteen-cell limitations.
% U11 (currency check of 2fe070fec): the GEANT4 block bound, DECISION-20261116-joseph-paper-batch.md item 3 (A6) wording, from fd/receipts-ll5d/mumax/run/node/mumax.json (422db056; commit e3a47f391), by mumax_values() in make_values_fived16.py.
% absence-ok: "remains untested" and "was not computed" are scoping statements fixed by the declaration (LL5D-DECLARATION.md :323-324, sec. 13 CANNOT SUPPORT) and the paper's approved section 7 sentence (DECISION-20261116-joseph-paper-batch.md item 2)

% Note round 1 (2026-11-08, draft for a PRD-referee read and Joseph): OUTCOME-20261106-1b-stage1-receipt-stop.md @ b551b05f :101-104 ("in the note and the report, a description of that rule and of the counts above"; the counts :27-30, the rule :72-74) and :106-114 (no coverage statement in either direction, no mechanism, nothing about 2D). Numbers:
% \lrval{stop-*}, checked against the record's lines 27 and 29 by make_values_landed.py. Factual correction 2026-11-06 (DECISION-20261106-joseph-diagnoses-jinst-corrections.md item 3): "is under way" was stale. OUTCOME-20261106-1b-stage1-receipt-stop.md (#832), "What this licenses": "the declared 5D coverage test (1B) did not run to a result": its stage-1 receipt stopped under a declared fail-closed rule (amendment 2 sec. 2b STOP, ending stage 1 under rule (iii)). NOT a coverage result in either direction. Caveat cut list MV02 (DECISION-20261102-joseph-caveat-cut-list.md; the director: "the staged method-coverage campaign is running toward a result"): FD-1B-STAGED-DECLARATION.md (stage 1 at G18_10a_02_11b and AR23_20i_00_000, the covariance remade at each truth); stage 0 PASS (FD-1B-STAGED-STAGE0-REPORT.md, rule K, kappa in its window in 16 of 16 cells); stage 1 stopped at its N1 guard before any toy (FD-1B-STAGED-STAGE1-STOP-REPORT.md) and runs on the redesign (DECISION-20261102-joseph-1b-redesign.md item 1) 620 round (the director, 2026-11-01, after #634): the stopped 1B test, FD-1B-DECLARATION.md sec. 1-2 (the arms at the two tunes; the matrix under test C_total + C_resp, no run-period block; the protocol fixed before any toy), FD-1B-DECLARATION.md sec. 3 (the draws table: lateral and normalization surrogates, C_model and C_resp not drawn), FD-1B-STAGE2-REPORT.md sec. 4.1 (gate D: rho_c = sqrt(E_d/M_d) outside [0.9, 1.1] in 11 of 16 cells, range 0.696-2.296; S1's stat, normalization and lateral blocks enter E_d unchanged; STOP before any coverage toy; gate 1 passed; "Not measured: why the ratios depart from 1"), FD-1B-DIAG-DECLARATION.md L1 (S1's diagonal in place of the toys' for every band in a set of at most K_MAX), FD-1B-DIAG-REPORT.md (#634: L1, no set of at most K_MAX = 3 of the 38 bands puts every cell in the window; NOT LOCALIZED; 1B stops); values_fived16.tex cov1b-* (gate-d.json e1838c0a, diag.json 6fd6669e). NOT a coverage result in either direction.
% Search: coverage of the sixteen-cell intervals; validated intervals of the sixteen-cell result; coverage at held-out truths (nd-unfolding/pv-ladder/readouts/cmodel_5d.json labels.coverage: "coverage at held-out truths was not validated"; FD-F0b-REPORT.md NOT CITABLE FOR coverage); the two-dimensional covariance not independently validated (paper sec:uq2d); fine-grid convergence (paper sec:conditions) @ 1f7f93da Note round n4: the body home of the abstract's sentence (note-n4/ledger.md, batch F).

\subsection{Dependence of the retained prior on estimator settings}
\label{sec:fived16-prr}
% U1 (currency check of 2fe070fec): the prior-retention study FD-PRR, from fd/receipts-prr/read/cluster/readout.json (9238f22b; commit 3163dfba1), by currency_values(), tier1_values(), prr_levels_values() and prr_s53_values() in make_values_fived16.py; wording as the paper\'s section 5.3 at 5bbb9d269, with DECISION-20261116-joseph-paper-batch.md item 4 for the settings sentences.
A study tests which settings of the \estMultiD{} change the fraction of a truth change that remains unrecovered. It uses noise-free pseudo-data with backgrounds at one split of the simulation and measures the retained fraction $\rho=-\sum_c b_c d_c/\sum_c d_c^2$ at each graded truth. Here $d_c=(X_c-P_c)/P_c$ is the truth's departure from the prior in cell $c$ and $b_c=(\hat X_c-X_c)/P_c$ is the bias, with $\hat X$, $X$ and $P$ the unfolded, true and prior cross sections. The sums run over cells where $|d_c|$ is at least \SI{1}{\percent}. Then $\rho=1$ if the result is the prior and $\rho=0$ if the departure is recovered; $\rho$ is one projection over the sixteen cells. The study changes one setting at a time: the iteration count, classifier capacity, training statistics or observables used by the classifiers. It also runs an exact binned IBU on the same pseudo-data. A level removes a truth's bias flag if the truth is no longer listed under the criterion of \S\ref{sec:fived16-fd} applied to the noise-free pseudo-data.

Along the prior-to-truth direction, IBU overshoots the truth: its retained fraction $\rho$ is $\fvval{prr-rho-ibu-gibuu}$ and $\fvval{prr-rho-ibu-q3}$, against $\fvval{prr-rho-e-gibuu}$ and $\fvval{prr-rho-e-q3}$ for the \estMultiD{}, and $\fvval{prr-rho-ibu200-gibuu}$ and $\fvval{prr-rho-ibu200-q3}$ at 200 iterations. The \estMultiD{}'s retained fractions are positive: it retains part of each truth's departure. The retained fraction is a directional summary and does not measure the total bias.

For each setting, $\rho$ at each tested level is compared with its value in the \estMultiD{} at each of the two truths flagged under the baseline covariance, and a change is reported when it differs by at least a threshold of \fvval{prr-floor}. A separate test asks whether IBU's retained fraction falls or rises between five and 200 iterations. The binned IBU is exact for the binned problem on its truth grid, while OmniFold uses event-level features, so a difference between the two combines the algorithm with the discretization.

No tested level of any setting lowers $\rho$ by as much as the threshold, and under the baseline covariance no tested level of a setting removes either truth's bias flag. Under $C_{\mathrm{rep}}$, the GiBUU hybrid truth stays flagged at every level, and the $q_3$ tilt is flagged only at the reduced-capacity and $(\pT,\pPar)$-only settings. Under the baseline covariance, IBU at five iterations removes the flag at the $q_3$ tilt ($p=\fvval{prr-ce-ibu5-q3-p}$) but not at the GiBUU hybrid truth, and newly flags \fvval{prr-ce-ibu5-nnew-word} of the other graded truths. Under $C_{\mathrm{rep}}$ it keeps the GiBUU hybrid truth's flag, leaves the $q_3$ tilt unflagged ($p=\fvval{prr-ibu5-q3-p}$), and newly flags \fvval{prr-ibu5-nnew-word} of the other graded truths. Reducing capacity to \fvval{prr-s53-c25-trees} trees and \fvval{prr-s53-c25-leaves} leaves, or restricting all three learners to $(\pT,\pPar)$, increased $\rho$ beyond the threshold at both truths. Dropping $W$ alone, or $W$ and $q_3$, changed $\rho$ by less than \fvval{prr-s53-w-dmax}. At both truths, IBU at five iterations has a lower retained fraction than the \estMultiD{}, by at least its threshold (\fvval{prr-ibu-bar}), and its retained fraction falls by at least a fixed margin between five and 200 iterations. At both, the retained prior is sensitive to the observable set. These results show that the retained prior is sensitive to the implementation and to some of its settings, but they identify no unique cause of it.

\subsection{Dependence on the unfolding algorithm}
\label{sec:fived16-ibu}
A comparison, pre-specified (Table~\ref{tab:test-index}), puts the result beside a binned D'Agostini iterative Bayesian unfolding (IBU) of the same data, with the same simulation, the same background template and the same weight-only systematic universes as the \estMultiD{}, at \lrval{rg1-K} iterations, integrated to the same sixteen cells. It forms $T=\Delta^{\top}C_{\Delta}^{-1}\Delta$ from the difference $\Delta$ of the two results and its covariance $C_{\Delta}$, which holds the statistical, universe and normalization blocks of the difference (not the model, lateral, $C_{\mathrm{resp}}$ or run-period blocks), and compares $T$ with bars calibrated on its own simulated null of exact agreement. The two results differ: $T=\lrval{rg1-T}$, against calibrated bars of \lrval{rg1-bar1} (\lrval{rg1-lev1}) and \lrval{rg1-bar2} (\lrval{rg1-lev2}). The reported cells therefore depend on the unfolding algorithm at the level of $T$; the largest per-cell differences and the ratio of the two methods' uncertainties are in the records document (Ref.~\cite{NoteRecords2026}). This dependence is a stated limitation of the result: no uncertainty block is added for it, and the reported cells and covariance are unchanged. A reading of differ means that a difference is detectable at this noise; it does not show that either method is correct.
% U14 (currency check of 2fe070fec): the per-cell differences from fd/receipts-ref-g1/a1-read/readout.json (4db170d3), printed_beside_R1 (fold_K5), over the reported sigma (sig-reported-*), by currency_values() in make_values_fived16.py; the paper's section 4.5 and limitation 4 at 5bbb9d269.

% Note round 1, N2 (2026-11-08, draft for a PRD-referee read and Joseph): DECISION-20261108-joseph-note-framing-and-estimator-dependence.md (#947, merged at af73565) item 3 (a stated limitation; no block; the 16 cells unchanged). OUTCOME-20261106-ref-g1-r1-differ.md @ b551b05f :40-41 (R1's licence, quoted), :46 (R2, no coverage statement), :47 (R3 descriptively; not "OmniFold adds discrimination", limitation 3 stands), :42-45 (no IBU switch, no block, no post hoc choice). Design: REF-G1-DECLARATION.md :25-26 (same data, simulation, template, universes), :82-83 (C_Delta's blocks), :130 (K = 5). Numbers: \lrval{rg1-*} from fd/receipts-ref-g1/a1-read/readout.json (4db170d3), checked against the record's lines 16-17 and 35-36 by make_values_landed.py; R3's note "(e)'s C_stat absorbs its estimator noise ... IBU has none" is the readout's.

